\section{Discussion, limitations, and future directions}
\label{sec:discussion}

\subsection{Limitations (Q4)}
We are deliberately conservative about what this study establishes.

\begin{enumerate}
  \item \textbf{The example is easy.} The model is two-dimensional with a simple
        linear reaction coordinate $\xib=x$ and a single dominant barrier. The
        orthogonal coordinate $y$ mixes quickly, so the conditional
        $Y\mid X=x$ equilibrates fast and the ABF feedback is forgiving. Effects
        seen here may not transfer to high-dimensional systems or nonlinear
        reaction coordinates.
  \item \textbf{ABF only is already a strong baseline.} Because plain ABF
        already converges well on this potential, the \emph{room} for FR to help
        is small; this is why the practical improvements are modest
        (\cref{sec:q1}) rather than transformative. A modest improvement over a
        strong baseline should not be oversold.
  \item \textbf{The oracle target is not a method.} The oracle uses $\Fref$, the
        unknown being computed, and is reported only as a diagnostic positive
        control (\cref{rem:oracle}). Its strong performance must be read as
        ``the FR mechanism has potential when the target direction is accurate'',
        never as a deployable result. All practical claims rest on the
        estimated target.
  \item \textbf{Conditional diagnostics are not configuration-resolved.} As in
        \cref{rem:cond_limit}, the conditional $Y\mid X$ output is keyed only by
        (method, target, seed). It supports the aggregate claim that successful
        FR does not corrupt the conditional, but cannot causally attribute
        good-versus-bad behaviour to specific hyperparameters.
  \item \textbf{The production estimator is a validated surrogate.} The GPU
        backend uses a binned-smoothed mean-force/marginal estimator and a
        vectorised fixed-$N$ birth--death. These pass the validation suite and
        agree with the kernel reference to within tolerance, but they are not
        mathematically identical to the original kernel notebook implementation
        (\texttt{ABF-FR.ipynb}); a residual discretisation difference exists by
        construction.
  \item \textbf{Harder regimes are untested.} The decisive question --- is FR
        ``icing on the cake'' (a mild accelerator when ABF already works) or
        ``snow in the winter'' (a rescue when ABF alone struggles)? --- cannot
        be answered from an easy example. It requires regimes that stress ABF:
        larger $\beta$ (higher barriers), fewer particles, shorter computational
        budgets, and slower orthogonal mixing (a stiff or metastable $Y\mid X$).
\end{enumerate}

\subsection{Failure modes}
The analysis points to two concrete ways FR can fail, both consistent with the
data. First, a \emph{noisy or biased target}: the estimated--oracle gap
(\cref{sec:q3}) shows that an inaccurate online target free energy caps the
attainable benefit, and a sufficiently wrong target would direct the
birth--death toward the wrong regions of $x$. Second, an \emph{overly aggressive
mechanism}: although the present configurations are gentle ($\sim10^{-5}$ event
fractions), large $\gamma$ or large event fractions could erode particle
diversity or perturb $Y\mid X=x$ faster than the dynamics can re-equilibrate it,
turning the coverage gain into a conditional bias. The uniform target is a third,
milder failure mode of \emph{target misspecification}: it works here only because
the converged biased marginal is near-uniform, and would misdirect FR once
$\Bt\not\approx F$.

\subsection{Future directions}
\begin{itemize}
  \item \textbf{Stress-test regimes.} Repeat the study at larger $\beta$, with
        fewer particles, shorter budgets, and a stiffer orthogonal coordinate,
        to determine whether FR is an accelerator or a rescue.
  \item \textbf{Better target estimation.} Since target quality is the practical
        bottleneck, improve the online estimate of the target free energy ---
        variance reduction, better EMA/smoothing, or uncertainty-aware targets
        --- to close the estimated--oracle gap.
  \item \textbf{Configuration-resolved conditional diagnostics.} Emit
        conditional $Y\mid X$ statistics per configuration so that good-versus-bad
        FR can be attributed mechanistically (\cref{rem:cond_limit}).
  \item \textbf{Adaptive FR strength.} Tie $\gamma$ and the burn-in to a
        data-driven signal (e.g.\ the measured imbalance or target uncertainty)
        rather than a fixed grid, to obtain the early-time speed of an aggressive
        rate without sacrificing final accuracy.
  \item \textbf{Higher-dimensional and nonlinear $\xib$.} Test multi-barrier and
        multi-dimensional reaction coordinates where the conditional sampling is
        genuinely hard.
\end{itemize}
